\section*{Chapitre \ref{ch:dv:managing-an-exotic-book} --- Gérer un livre exotique}

\begin{solution}{exo:dv:managing-an-exotic-book:1}
Sa valeur dépend fortement de la corrélation des deux indices sur trois ans, grandeur qu'aucun marché ne cote~: une donnée non observable, de niveau~3. Elle pourrait passer au
niveau 2 si la corrélation devenait observable pendant presque toute la durée restant à courir du titre (\omterm{def:dv:multi-asset-options:corrswap}{swaps de corrélation} ou paniers cotés sur la paire), ou si l'effet de
la corrélation sur la valeur devenait négligeable.
\end{solution}

\begin{solution}{exo:dv:managing-an-exotic-book:2}
La banque est débitrice du titre. Plus la valeur du titre est élevée, plus la valeur comptable de la banque est faible~: le point prudent, c'est-à-dire la valeur à laquelle elle a 90~\% de chances de sortir ou mieux, se situe à
la borne haute des valeurs plausibles du titre. Avec quatre modèles, le point à 90~\% se situe sous la première position de tracé, et c'est la valeur la plus élevée~: 98.62.
\end{solution}

\begin{solution}{exo:dv:managing-an-exotic-book:3}
$0.621+0.621+0.433+0.325=2.000$ millions. Différer le montant net de bid--offer, sans rien constater à l'initiation, reporte $2.000-0.621=1.379$ million.
\end{solution}

\begin{solution}{exo:dv:managing-an-exotic-book:4}
Pas par construction~: la réserve bid--offer chiffre la fermeture du \omterm{def:dv:greeks-and-the-hedging-pnl:greeks}{vega} connu, tandis que la \omterm{def:dv:managing-an-exotic-book:model}{réserve de modèle} valorise sans savoir quelle volatilité est la bonne. Elles coïncident
parce que le modèle alternatif se situe ici à un demi-point de volatilité, et que la fourchette vaut un demi-point. Un autre ensemble de modèles, ou une autre fourchette, les sépare.
La norme européenne demande que le risque de modèle n'inclue pas ce que les réserves de prix de marché couvrent déjà, si bien que le contrôle doit figurer dans la documentation de la réserve.
\end{solution}

\begin{solution}{exo:dv:managing-an-exotic-book:5}
Le titre rembourse selon le moins bon des deux indices. Avec une corrélation plus élevée, les indices évoluent ensemble, le moins bon se situe moins loin sous l'autre, et le knock-in et les
coupons manqués sont moins probables~: le titre vaut davantage pour son détenteur.
\end{solution}

\begin{solution}{exo:dv:managing-an-exotic-book:6}
La banque est longue de \omterm{def:dv:greeks-and-the-hedging-pnl:greeks}{vega}~: cinq points de volatilité en moins coûtent à eux seuls 5.89 million. Un krach atteint le knock-in, où le put long de la banque rapporte et compense la perte de volatilité
(7.02 million à $-30\%$). Une baisse modérée rapporte peu à volatilité inchangée (1.31 million à $-10\%$), et à volatilité plus faible le knock-in est encore moins probable, si bien que
rien ne compense la perte de \omterm{def:dv:greeks-and-the-hedging-pnl:greeks}{vega}~: $-6.31$ million à $-10\%$ avec cinq points de volatilité en moins.
\end{solution}

\begin{solution}{exo:dv:managing-an-exotic-book:7}
Pour une corrélation comprise entre 0.2 et 0.8, le titre vaut 97.33 à 99.12~; le point à 90~\% vaut 98.93 et la réserve 0.935. Le montant différé devient $0.621+0.935=1.556$ et le
\omterm{def:dv:managing-an-exotic-book:dayone}{P\&L du premier jour} constaté $2.000-0.621-1.556=-0.177$ million~: à cet intervalle, la vente au pair ne couvre pas ses réserves.
\end{solution}

\begin{solution}{exo:dv:managing-an-exotic-book:8}
Additionner chaque réserve à son point à 90~\% suppose que chaque position soit liquidée à son pire niveau au même moment. Ce n'est pas prudent par construction~: c'est mal calibré, et cela rend
la réserve inutile comme mesure de ce que le livre pourrait perdre. L'agrégation de la norme de valorisation prudente (la moitié de la somme selon sa première méthode) existe pour cette
raison. À l'inverse, les positions compensatoires doivent être nettées avant d'être réservées, faute de quoi le calcul compte deux fois une même exposition.
\end{solution}

\begin{solution}{pb:dv:managing-an-exotic-book:1}
\textbf{1.} 7.95~\% par an. La banque est longue du put de knock-in des investisseurs (longue de \omterm{def:dv:greeks-and-the-hedging-pnl:greeks}{vega}, longue de skew) et courte en corrélation.
\textbf{2.} Observable~: les indices, les taux, les dividendes, les volatilités à la monnaie. Non observables~: la corrélation sur trois ans et le choix de la volatilité pour le knock-in.
\textbf{3.} Niveau 3. L'écart du premier jour ne peut être constaté que si la fair value est attestée par des cours de niveau 1 ou par des données uniquement observables~; sinon, il est différé.
\textbf{4.} \omterm{def:dv:greeks-and-the-hedging-pnl:greeks}{Vegas} $+0.533$ et $+0.693$ par point et par tranche de 100~; deltas $-0.143$ et $-0.279$ par 1~\% avant couverture (la banque est courte du titre).
\textbf{5.} Les deltas, et une partie du \omterm{def:dv:greeks-and-the-hedging-pnl:greeks}{vega} avec des options cotées~; ni la corrélation, ni le skew de long terme.
\textbf{6.} 0.266 et 0.346 pour les \omterm{def:dv:greeks-and-the-hedging-pnl:greeks}{vegas} et 0.008 pour les deltas~: 0.621 million (après arrondi).
\textbf{7.} 98.00, 94.33, 98.62, 97.98~; \omterm{def:dv:managing-an-exotic-book:model}{réserve de modèle} de 0.621 million.
\textbf{8.} 97.64 à 98.50~; le point à 90~\% vaut 98.43~; réserve de 0.433 million.
\textbf{9.} La moitié de $0.621+0.621+0.433$~: 0.838 million.
\textbf{10.} 69 jours, avec une hypothèse de 100\,000 par point et par jour et un dixième du volume~: plus de dix, donc un ajustement de concentration est dû.
\textbf{11.} Passé en charge 0.621 (bid--offer), différé 1.054, constaté 0.325 million.
\textbf{12.} La corrélation devient observable dans l'intervalle 0.47--0.53 et il reste deux ans au titre~: la réserve de corrélation tombe à 0.071 et la \omterm{def:dv:managing-an-exotic-book:model}{réserve de modèle} à 0.559~; 0.424
million est repris.
\textbf{13.} 0.362 million.
\textbf{14.} Avec les indices à 80~\%, le titre vaut 86.99 et se rapproche de sa barrière~: la \omterm{def:dv:managing-an-exotic-book:model}{réserve de modèle} monte à 0.649 et la réserve de corrélation vaut 0.202, si bien que seuls 0.203
million est repris.
\textbf{15.} La couverture devrait protéger une baisse modérée accompagnée d'une baisse de volatilité~: vendre moins de couverture à la hausse, ou conserver un peu de \omterm{def:dv:greeks-and-the-hedging-pnl:greeks}{vega} à court terme à mesure que les indices baissent, dans les limites fixées.
\textbf{16.} Lent dans le temps, prudent dans le montant~: le profit ne revient qu'une fois l'incertitude levée, et si la levée tourne mal, une partie ne revient jamais.
\textbf{17.} Elle réduit la rémunération des marges qui reposent sur des données non observables, ce qui est son objet, et elle pousse les desks à négocier des structures dont les données peuvent être observées.
\textbf{18.} Une fonction indépendante de contrôle de la valorisation, et non le desk~: le desk propose, le contrôle de la valorisation décide et documente.
\textbf{19.} Sur une marge de 2 millions pour 100 millions de titres en notionnel, 0.325 million est constaté à l'initiation (0.621 passé en charge, 1.054 différé)~; 0.424 million est repris au bout d'un an,
dont 0.362 lorsque la corrélation devient observable.
\textbf{20.} Valoriser une position au prix auquel la banque pourrait la liquider, et non à ce que dit son modèle.
\end{solution}

\begin{solution}{iq:dv:managing-an-exotic-book:1}
L'écart à l'initiation entre le prix de la transaction et la fair value. Selon IFRS 9, il n'est constaté immédiatement que si la fair value est attestée par un prix de niveau 1 ou
par une valorisation fondée uniquement sur des données observables~; sinon il est différé, puis repris au rythme des facteurs que les acteurs de marché intégreraient dans leur prix, dont le temps.
\iqlookfor{observabilité et report.}
\end{solution}

\begin{solution}{iq:dv:managing-an-exotic-book:2}
Valoriser l'exotique selon un ensemble de modèles qui reproduisent les mêmes vanilles (\omterm{def:dv:local-volatility:model}{volatilité locale}, stochastique, locale-stochastique~; calibrations alternatives), en retenir l'intervalle, et réserver la
distance entre la valeur comptabilisée et un point prudent (le point à 90~\% de la norme européenne). Documenter l'ensemble, le revoir à mesure que les marchés évoluent, et éviter le double comptage avec les
réserves de prix de marché.
\iqlookfor{ensemble de modèles, point prudent, indépendance.}
\end{solution}

\begin{solution}{iq:dv:managing-an-exotic-book:3}
Position longue en volatilité sur les échéances lointaines, longue en skew, courte en corrélation (worst-of), sensible aux dividendes, et gamma important près des barrières à mesure qu'elles approchent. Elles sont difficiles à couvrir parce que peu
d'intervenants de marché prennent la contrepartie, et les expositions sont partagées par chaque émetteur.
\iqlookfor{les positions de \omterm{def:dv:greeks-and-the-hedging-pnl:greeks}{vega} et de corrélation et leur marché à sens unique.}
\end{solution}

\begin{solution}{iq:dv:managing-an-exotic-book:4}
Une grille de variations d'indices et de volatilité, avec des chocs de corrélation et de dividendes, une revalorisation complète et les couvertures en place~; à quoi s'ajoutent les scénarios de trajectoire dans lesquels les indices approchent les
barrières. La situation redoutée ici est une baisse modérée accompagnée d'une baisse de volatilité, où les couvertures perdent et le \omterm{def:dv:greeks-and-the-hedging-pnl:greeks}{vega} long aussi, et non le krach.
\iqlookfor{revalorisation complète, chocs conjoints, le pire cas non évident.}
\end{solution}

\begin{solution}{iq:dv:managing-an-exotic-book:5}
L'instantané de données de marché, le modèle et sa version, l'ensemble de modèles et les intervalles plausibles avec leur justification, les positions, la version du code et les valeurs calculées à chaque
étape, de sorte que les mêmes entrées reproduisent la même réserve.
\iqlookfor{instantané, versions, intervalles et leur justification.}
\end{solution}

\begin{solution}{iq:dv:managing-an-exotic-book:6}
La réserve bid--offer de paramètre associée se réduit à la fourchette du marché, et le \omterm{def:dv:managing-an-exotic-book:dayone}{P\&L du premier jour} différé qui lui est lié est repris~; la valorisation se déplace vers la valeur observée, qui peut être un gain
ou une perte par rapport au mid comptabilisé.
\iqlookfor{reprise du report, et effet de valorisation au prix de marché.}
\end{solution}
